\chapter*{Notations et abréviations}
\addcontentsline{toc}{chapter}{Notations et abréviations}
\markboth{Notations et abréviations}{Notations et abréviations}

\section*{Abréviations du procédé}

\begin{longtable}{@{}ll@{}}
\toprule
\textbf{Sigle} & \textbf{Signification} \\
\midrule
\endhead
P29 & Acide phosphorique à \SI{29}{\percent} de \PdeuxOcinq{} (acide « faible ») \\
P54 & Acide phosphorique à \SI{54}{\percent} de \PdeuxOcinq{} (acide « fort ») \\
NCL & Non clarifié --- acide 54 brut sortant de la concentration \\
CL & Clarifié --- acide 54 débarrassé de ses impuretés solides \\
DEC & Décadmié --- acide dont le cadmium a été retiré \\
CoC & Cocristallisé --- acide ultra-purifié par cristallisation \\
DEC\_CL & Décadmié puis clarifié \\
IR11 & Bac central recevant le CoC et le DEC\_CL \\
IR12 & Bac central recevant l'acide clarifié ordinaire \\
ARP & Acide de retour phosphorique (raffinat d'EMAPHOS) \\
MILP & \emph{Mixed Integer Linear Programming} --- programmation linéaire mixte \\
\bottomrule
\end{longtable}

\section*{Ensembles}

\begin{longtable}{@{}ll@{}}
\toprule
\textbf{Symbole} & \textbf{Définition} \\
\midrule
\endhead
$\Lvingtneuf$ & Lignes de production d'acide 29 (6 éléments) \\
$\Lcinqquatre$ & Lignes de concentration (5 éléments) \\
$\Lvingtneufetoile$ & Lignes d'acide 29 disposant d'une concentration associée \\
$\Ech$ & Échelons de concentration (24 éléments) \\
$\Ech_m$ & Échelons de la ligne $m$ \\
$\Ech^{\mathrm{coc}}_m$ & Échelons de $m$ affectés à la cocristallisation \\
$\Conso$ & Consommateurs (9 éléments) \\
$\Prod$ & Produits engrais \\
$\Acides$ & Types d'acide (6 éléments) \\
$\sigma$ & Application ligne 29 $\to$ ligne 54 \\
$\Gamma$ & Liaisons interzones autorisées (14 arcs) \\
$\Phi_a$ & Raccordements ligne $\to$ consommateur pour l'acide $a$ \\
\bottomrule
\end{longtable}

\section*{Paramètres}

\begin{longtable}{@{}lll@{}}
\toprule
\textbf{Symbole} & \textbf{Définition} & \textbf{Unité} \\
\midrule
\endhead
$P_\ell$ & Production d'acide 29 de la ligne $\ell$ & \tP/j \\
$C_e$ & Capacité de l'échelon $e$ à plein régime & \tP/j \\
$h_e$ & Heures de marche de l'échelon $e$ & h \\
$\kappa_e$ & Production de l'échelon $e$ & \tP/j \\
$\Pi_m$ & Production totale de la ligne $m$ & \tP/j \\
$\Picoc_m$ & Part produite par les échelons cocristallisants & \tP/j \\
$\Pincl_m$ & Part produite par les autres échelons & \tP/j \\
$G_m$ & CoC produit (systématique) & \tP/j \\
$B^{\mathrm{coc}}_m$ & Boue de cocristallisation & \tP/j \\
$\etaconc$ & Rendement de concentration $= 1{,}00$ & --- \\
$\etacl$ & Rendement de clarification $= 0{,}90$ & --- \\
$\etacoc$ & Rendement de cocristallisation $= 0{,}80$ & --- \\
$\Lambda$ & Capacité entrante des décanteurs $= 1000$ & \tP/j \\
$\mathcal{D}$ & Niveaux de décadmiation $= \{0, 750, 1500\}$ & \tP/j \\
$\tau^{\min}, \tau^{\max}$ & Bornes d'un transfert interzone & \tP \\
$\alpha$ & Part minimale de DEC\_CL dans IR11 & --- \\
$\epsilon$ & Tolérance sur la charge de concentration & --- \\
$R_{k,a}$ & Besoin consolidé du consommateur $k$ en acide $a$ & \tP \\
$\beta_\ell$ & Part du retour EMAPHOS vers la ligne $\ell$ & --- \\
$Z^{\min}, Z^{\max}$ & Bornes de la bande de sécurité d'un stock & \tP \\
\bottomrule
\end{longtable}

\section*{Variables de décision}

\begin{longtable}{@{}lll@{}}
\toprule
\textbf{Symbole} & \textbf{Définition} & \textbf{Nature} \\
\midrule
\endhead
$w_{\ell,v}$ & Choix du niveau $v$ de décadmiation sur $\ell$ & binaire \\
$d_\ell$ & Quantité décadmiée sur $\ell$ & dérivée \\
$\xstd_\ell,\ \xdec_\ell$ & Acide 29 envoyé en concentration & continue \\
$t_{\ell\ell'}$ & Transfert interzone & semi-continue \\
$b_{\ell\ell'}$ & Activation du transfert & binaire \\
$\uncl_m$ & NCL envoyé en clarification & continue \\
$y$ & Livraisons aux consommateurs & continue \\
$S$ & Stocks finaux & continue \\
$\rho_{k,a}$ & Demande non satisfaite & continue \\
$\nu^{+}, \nu^{-}$ & Écarts aux bandes de stock & continue \\
$\varepsilon^{+}, \varepsilon^{-}$ & Écarts à la charge planifiée & continue \\
$g^{+}, g^{-}$ & Écarts des stocks finaux à leur cible & continue \\
\bottomrule
\end{longtable}

\cleardoublepage
